% Generated by `taskgen course examples` from dlp.index_calculus (seed 3).
% Do not edit: change examples.yaml and regenerate.
\begin{example}[метод исчисления индексов целиком]\label{ex:ic-full}
Решить задачу дискретного логарифмирования методом Исчисления индексов: \[6^x \equiv 62\ (\textrm{mod}\ 83).\] Обосновать выбор значения параметра $B$. Систему необходимо выписать и решить в явном виде.

\par\noindent\textit{Решение.} $g = 6$ --- первообразный корень, логарифмы берутся по модулю $p - 1 = 82 = 2 \cdot 41$. Выбор $B$: $L(p) = e^{\sqrt{\ln p \ln\ln p}} \approx 12.97$, $B \approx L(p)^{1/\sqrt 2} \approx 6.12$; берём $B = 6$, факторная база $\{2, 3, 5\}$, неизвестные $L_\ell = \log_g \ell$.\par Вычисляем $g^k \bmod p$, $k = 1, 2, \dots$ и оставляем $B$-гладкие значения (не гладкие: $g^{4} \equiv 51$, $g^{5} \equiv 57$). Достаточная для однозначного решения система:\[\begin{array}{rcl}g^{1} \equiv 6 = 2 \cdot 3 & \Rightarrow & 1 \equiv L_{2} + L_{3} \\ g^{3} \equiv 50 = 2 \cdot 5^{2} & \Rightarrow & 3 \equiv L_{2} + 2L_{5} \\ g^{6} \equiv 10 = 2 \cdot 5 & \Rightarrow & 6 \equiv L_{2} + L_{5} \\ \end{array}\pmod{82}\]Решаем систему по модулю каждого примарного делителя $p - 1$ (метод Гаусса; ведущие коэффициенты обратимы), ступенчатая форма (в строке --- коэффициенты при $L_{2}$, $L_{3}$, $L_{5}$ и правая часть):\par\noindent по модулю $2$: $(1, 1, 0 \mid 1)$, $(0, 1, 0 \mid 0)$, $(0, 0, 1 \mid 1)$ $\Rightarrow$ $L_{2} \equiv 1$, $L_{3} \equiv 0$, $L_{5} \equiv 1$;\par\noindent по модулю $41$: $(1, 1, 0 \mid 1)$, $(0, 1, 39 \mid 39)$, $(0, 0, 1 \mid 38)$ $\Rightarrow$ $L_{2} \equiv 9$, $L_{3} \equiv 33$, $L_{5} \equiv 38$;\par\noindent по КТО: $L_{2} = 9$, $L_{3} = 74$, $L_{5} = 79$ (проверка: $6^{9} \equiv 2$, $6^{74} \equiv 3$, $6^{79} \equiv 5$).\par Ищем $s$ с $B$-гладким $h g^{s}$: $s = 1$, $62 \cdot 6^{1} \equiv 40 = 2^{3} \cdot 5$, откуда $x \equiv 3L_{2} + L_{5} - 1 \equiv 23\ (\textrm{mod}\ 82)$.\[x = 23\ (\textrm{mod}\ 82).\]
\end{example}
